% !PRINTED 103
\weberpage{103}
%<<CONT
by \S.\,22 one obtains for the indices of the numbers in $\mathfrak{A}$ the two congruences:
\weq{3}{\begin{aligned}
e_{1,1}\alpha_1 + e_{1,2}\alpha_2 + \cdots + e_{1,\mu}\alpha_\mu &\equiv 0\\
e_{2,1}\alpha_1 + e_{2,2}\alpha_2 + \cdots + e_{2,\mu}\alpha_\mu &\equiv 0
\end{aligned}\ \ (\mathrm{mod}\ 2).}
To determine the number of possible groups $\mathfrak{A}$, one should bear in mind that the relations (3) can be replaced, without changing their content, by two mutually independent linear combinations. Accordingly one may take one of the pairs of numbers in the matrix (2), e.\,g. $e_{1,1}, e_{2,1} = 1, 0$, and thus put in place of (2):
\[
\begin{array}{cccc}
1, & e_{1,2}, & \ldots & e_{1,\mu}\\
0, & e_{2,2}, & \ldots & e_{2,\mu}.
\end{array}
\]
Now for the $e_{2,2} \ldots e_{2,\mu}$ one can put any combination of the numbers 0, 1, with the exception of the one in which all the $e_{2,h} = 0$. To every $e_{2,h} = 0$ the corresponding $e_{1,h} = 1$ must hold, while to an $e_{2,h} = 1$ there may correspond an $e_{1,h} = 0$ or $= 1$. The number of combinations arising in this way is:
\[
1 + (\mu - 1)2 + \frac{(\mu - 1)(\mu - 2)}{1.2} 2^2 + \cdots + 2^{\mu - 1} - 1 = 3^{\mu - 1} - 1.
\]
But now the second of the equations (3) can, without changing its meaning, also be added to the first, so that every two of the combinations obtained yield the same group. There thus remain
\weq{4}{\frac{3^{\mu - 1} - 1}{2}}
distinct combinations of the $e_{h,k}$, which also, as is easily seen, lead to distinct groups $\mathfrak{A}$. For suppose that the one time one has taken $e_{1,2}, e_{2,2} = 1, 0$, the other time $= 0, 1$; then $\alpha_1 = 0, \alpha_2 = 1, \alpha_3 = 0, \ldots \alpha_\mu = 0$ satisfies the first assumption, but does not satisfy the second. % [sic: Weber prints e_{h,k}; the defining matrix and congruences consistently use e_{k,h}.]

The number that we have just determined is exhaustive only when $m$ is not divisible by 8, because it is thereby assumed that we never have $e_{1,h}, e_{2,h} = 0, 0$. But if $m$ is divisible by 8, then this combination of values can occur in one pair, and we must therefore still add the cases which, through
\[
\begin{array}{ccccc}
1 & 0 & e_{1,3} & \ldots & e_{1,\mu}\\
0 & 0 & e_{2,3} & \ldots & e_{2,\mu}
\end{array}
\]
%CONT>>
